% Hand-written sentences that depend on results; every number is a macro from v19_numbers.tex (Research/code/v19_tables.py).
\newcommand{\nPointerSentence}{by its certificates, fixed forms lose $4.00$ and $8.00$, the maximum that Theorem~\ref{thm:universal} allows, two rounds lose $1.80$ and between $3.32$ and $3.33$, against the upper bounds $2$ and $4$ of part (a), and $b-1$ rounds lose $1.80$ at both caps (Appendix~\ref{app:pointerexact}).}
\newcommand{\nOrderSentence}{an illustrative flip of the last answer with probability 10\% biases the scores by \nPosBranch--\nPosFit{} standard errors at $n=1000$ (Appendix~\ref{app:measurement}).}
\newcommand{\nEmAppendixText}{At $n=1000$ the working model on the fitted-form records has the smallest RMSE of the four estimators in most settings, \nEmBrLo--\nEmBrHi{} times that of branching: the model shares strength across paths that the design-based scores treat separately. Its bias is between \nEmBiasLo{} and \nEmBiasHi{} branching standard deviations and its coverage \nEmCovLo--\nEmCovHi{} (below 0.925 in \nEmBelow{} of \nEmSettings{} settings). The bias is a property of the model, not of the sample size: the accuracy of a depth-$b$ tree depends on interactions of up to $b$ items and the label, which a two-way model does not represent. At $n=\nEmBigN$ the same bias is up to \nEmBigBiasHi{} of the estimator\textquotesingle s own standard deviations, coverage falls to \nEmBigCovLo--\nEmBigCovHi{} (below 0.90 in \nEmBigBelow{} of \nEmBigSettings{} settings), and the RMSE is \nEmBigRmseLo--\nEmBigRmseHi{} times that of branching and above that of the design-based form score in \nEmBigWorse{} settings; the design-based scores keep coverage \nEmBigFixCovLo--\nEmBigFixCovHi. A richer model would reduce the bias and raise the variance; with all interactions it is the design-based estimator. On balanced random forms the model-based estimate is unbiased to a similar degree but \nEmUnifLo--\nEmUnifHi{} times less accurate at $n=1000$ than the design-based score on optimized forms, so the choice of forms matters for model-based analysis as well.}
\newcommand{\nDiscussionSentence}{On the trained rules we examined, optimized forms needed between $1.0$ and \nMaxAll{} times the records of branching without a model, one adaptive round recovered at least \nTwoMinPct\% of the log-gap wherever its optimum was computed, and under a pairwise model the price of pre-commitment fell to \nModelTwoLo--\nModelTwoHi.}
\newcommand{\nAbstractData}{For screening trees and forest updates trained on survey and clinical data, optimized forms need between $1.0$ and \nMaxAll{} times the records of the adaptive optimum, one further round of adaptivity recovers most of the gap, and a pairwise working model lowers the price to at most \nModelAllHi, at the cost of a bias that does not shrink with the audit.}
\newcommand{\nContribData}{On screening trees for several health outcomes and on forest updates, optimized forms need $1.0$--\nMaxAll{} times the records of branching without a model; much of what a naive comparison credits to adaptivity is recovered by optimizing the forms, and a second round recovers most of the rest. Under a pairwise working model the price falls to \nModelTwoLo--\nModelAllHi, but the model's estimate has a bias that does not shrink with the audit, and its likelihood can be ill-posed on designs not chosen for it.}
\newcommand{\nAuditParagraph}{Audits of 500 to 32,000 respondents on \nCurveTrees{} trees, collected on fitted forms, on forms optimized for the pairwise model and by branching, show the same orders in finite samples (Figure~\ref{fig:rounds}b, Table~\ref{tab:curve}). The mean squared error of forms relative to branching is \nAudDesignLo--\nAudDesignHi{} for the design-based scores (medians over trees at each size), \nAudThreeFormsLo--\nAudThreeFormsHi{} for maximum likelihood under the three-way model, and \nAudTwoFormsSurveyLo--\nAudTwoFormsSurveyHi{} under the pairwise model up to 8,000 respondents, where it also matches the design-based mean under branching (\nAudTwoVsBranchSurveyLo--\nAudTwoVsBranchSurveyHi); the pairwise estimate is an early-stopped EM iterate, which is what remains defined where the likelihood is ill-posed. At 32,000 respondents the bias dominates: the pairwise fit is less accurate than the design-based form score for \nCurveTwoWorsed{} of the \nCurveTrees{} trees, and its intervals cover \nEmBigCovLo--\nEmBigCovHi{} (Table~\ref{tab:em}). A cross-fitted score whose tables come from a pairwise fit to the audit's own records stays unbiased, with coverage \nAipwCovLo--\nAipwCovHi{} over trees and sizes, but has the mean squared error of the reference tables (\nAipwMseLo--\nAipwMseHi) and does not inherit the model's small-sample gain (Table~\ref{tab:aipw}).}
\newcommand{\nAuditAppendixText}{For the design-based scores the ratio of mean squared errors stays near the population ratio at every size (medians \nAudDesignLo--\nAudDesignHi). For the three-way model the ratio between forms and branching records is \nAudThreeFormsLo--\nAudThreeFormsHi{} on the fitted forms and \nAudThreeMformsLo--\nAudThreeMformsHi{} on the forms optimized for the pairwise model; for the pairwise model it is \nAudTwoFormsSurveyLo--\nAudTwoFormsSurveyHi{} and \nAudTwoMformsSurveyLo--\nAudTwoMformsSurveyHi{} up to 8,000 respondents (at most \nAudTwoFormsSurveyMax{} and \nAudTwoMformsSurveyMax{} for a single tree). These orders of magnitude agree with the population constants of Table~\ref{tab:model}, although the target laws do not satisfy the models. Branching here is the rule's own tree, not the branching design that is optimal within the model, which Table~\ref{tab:model} puts a few percent below the optimal forms. At large sizes the comparison is governed by the bias of the model, on either kind of record: at 32,000 respondents the pairwise fit has a median absolute bias of \nAudBiasFormsBig{} of its standard deviation on forms (\nAudBiasFormsSmall{} at 1,000) and \nAudBiasBranchBig{} on branching records (\nAudBiasBranchSmall{} at 1,000), the pairwise fit to branching records is less accurate than the plain mean of $t$ for \nAudBranchTwoWorsec{} and \nAudBranchTwoWorsed{} of the \nCurveTrees{} trees at 16,000 and 32,000, and the ratio between the two designs ranges from \nAudTwoFormsMinAll{} to \nAudTwoFormsMaxAll{} over trees. Against the design-based form score on the same fitted-form records, pairwise maximum likelihood has the smaller mean squared error for every tree up to $n=\nCurveLastAllBetter$ and the larger for \nCurveTwoWorsec{} and \nCurveTwoWorsed{} of the \nCurveTrees{} trees at 16,000 and 32,000 (RMSE up to \nCurveTwoMaxc{} and \nCurveTwoMaxd{} times larger); which trees these are cannot be read from the capped records. A model-assisted score, which uses a working model only to choose the tables of a pointwise-unbiased score, as in the smoothed pilot of Appendix~\ref{app:pilot}, cannot have a smaller variance than the oracle tables for the same form distribution; the advantage of maximum likelihood at small $n$ is bought with the model's bias.}
\newcommand{\nForestScaleSentence}{On twenty Spambase features the picture is the same at $b=3$: forms need \nFScaleThreeLo--\nFScaleThreeHi{} times the records of branching and two rounds at most \nFScaleThreeTwoHi{} (Table~\ref{tab:forestscale}).}
\newcommand{\nForestScaleAppendixText}{At $b=3$ column generation over branching interviews closed within its budget on all three splits, so the ratios \nFScaleThreeLo--\nFScaleThreeHi{} are exact to the certified gap; they lie in the range of the twelve-feature tasks. The two-round programs did not close in 15 minutes, and their certified intervals put two rounds between \nFScaleThreeTwoLo{} and \nFScaleThreeTwoHi{} times the records of branching. At $b=4$ (4,845 forms) the fixed optimum was certified in seven minutes, but 25 minutes of column generation over branching interviews left only the trivial lower bound for the one split that was run, so its interval is uninformative; the other two splits were not run. Exact adaptive designs at $m=20$ and $b\ge4$ are beyond this implementation.}
\newcommand{\nDependenceAppendixText}{Along the strongly coupled path the pairwise constants of both classes rise from $1/3$ towards their values without a model, and their ratio rises from $1.00$ at $\beta=0$ to \nDepRatioLast{} at $\beta=\nDepBetaMax$, above the limit $2^{b-1}/b=1.33$ guaranteed by the theorem and below the factor 4 of part (a). At moderate coupling the gap is still small (\nDepRatioOne{} at $\beta=1$ and \nDepRatioTwo{} at $\beta=2$); it opens where the auxiliary features become nearly deterministic functions of the address and the data. At the pairwise fits of the noisy address populations the ratio stays at \nDepFitTwoMax{} at every accuracy, although the ratio without a model is $2^{b-1}$ at each of these populations: the association of the label with the data bits, which is what accuracy adds, does not bring the address structure into the pairwise tangent space. At the three-way fits the ratio is \nDepFitThreeLo--\nDepFitThreeHi.}
\newcommand{\nAipwAppendixText}{The cross-fitted estimator has the mean squared error of the reference-table score to within a few percent (median ratio \nAipwMseLo--\nAipwMseHi{} over the three sizes, at most \nAipwMseMax{} for a single tree), a bias within Monte Carlo error ($|$bias$|$/SD at most \nAipwBiasMax, against \nAipwBiasRefMax{} for the reference-table score on the same records), and coverage between \nAipwCovLo{} and \nAipwCovHi{} over trees and sizes (\nAipwCovRefLo--\nAipwCovRefHi{} for the reference-table score; the low values are for the same tree and size). A design-unbiased analysis therefore needs no external reference. It does not capture the small-sample advantage of pairwise maximum likelihood, as expected: for given form probabilities its variance cannot fall below that of the oracle tables, and the advantage of maximum likelihood is bought with the model's bias.}
\newcommand{\nPracticeSentence}{For an audit of a known rule under a per-record cap, the results suggest the following. If the rule uses at most $b$ features, ask them all; nothing is lost. Otherwise, if the interview can branch, administer the rule: its accuracy is then estimated by a plain mean with no design to compute. If questions must be pre-committed, do not randomize uniformly; solve for the optimal forms on a reference sample or on a pilot that is charged to the audit (a few percent to about a fifth more records than the oracle here), and if one batch of follow-up questions is possible, use it, since two rounds were within \nTwoBoundMax{} of branching in every screening run. A pairwise working model makes pre-committed forms nearly as informative as branching records for audits of a few thousand records, but its bias does not shrink and its intervals under-cover in large audits, so claims that must not rest on a model should use the design-based scores and pay the constants of Table~\ref{tab:screening}.}
\newcommand{\nRulesSentence}{Trees for four other BRFSS outcomes (heart disease, stroke, high blood pressure, difficulty walking; 21 variables) and for Spambase (57 features), with forms on the rule's items, give \nRulesLo--\nRulesHi{} in the \nRulesOk{} runs with a gain that were computed (median \nRulesMed); bootstrap resamples of the target records move a ratio by at most $\nRulesBootBrfss\%$ on BRFSS and $\nRulesBootSpam\%$ on the \nRulesSpamN{} Spambase target records (5th to 95th percentile), so the spread is between rules, not sampling error of the target law (Table~\ref{tab:rulesmain}).}
\newcommand{\nMisspecSentence}{The accuracy of a depth-$b$ tree involves interactions of up to $b$ items and the label, so the target laws do not satisfy the pairwise model, and the limit of its maximum-likelihood estimate depends on the design. At the target laws themselves we maximize the expected log-likelihood of the pairwise model under each design from \nMisStarts{} starts and compute the sandwich variance at the maximizer \citep{white1982,golden2019} (\nMisRuns{} runs; Table~\ref{tab:misspec}). Whether the estimator is well-defined depends on the design: all starts reach one interior maximizer in \nMisWellMF{} of the \nMisRuns{} runs for the forms that are optimal for the model, in \nMisWellMB{} for the model-optimal branching design, in \nMisWellF{} for the optimal design-based forms and in \nMisWellT{} for the rule's own tree, for which no interior maximizer is found in \nMisNoneT{} runs. On the model-optimal forms the limit misses the true accuracy by at most \nMisBiasPtsMf{} points, so the squared bias equals the variance at a median of \nMisCrossMf{} records, and where both model-optimal designs have an interior maximizer the sandwich variance of the forms is \nMisVarMmLo--\nMisVarMmHi{} times that of the branching design (median \nMisVarMmMed); the model's own variance formula understates the sandwich variance by up to a factor \nMisSandModelHi. Under misspecification the model-based analysis is thus best defined on pre-committed forms chosen for the model, and adaptivity adds little to it.}
\newcommand{\nRulesAppendixText}{The \nRulesAll{} runs comprise \nRulesOk{} with a gain, \nRulesTrivial{} in which the rule depends on at most $b$ items, and \nRulesSkipped{} that were not computed (\nRulesTooMany{} rules depend on more than 13 items, and \nRulesUnfinished{} runs, at caps five and six, were not finished within the time allowed). The ratios lie in the range of the diabetes trees of Table~\ref{tab:screening}, with one larger value (\nRulesHi). Resampling the target records with the tree fixed moves a ratio by at most \nRulesBootBrfss\% of its value on BRFSS, where a target half has more than $10^5$ records, and by at most \nRulesBootSpam\% on Spambase (\nRulesSpamN{} target records; 5th to 95th percentile of \nRulesBootReps{} resamples, \nRulesBootN{} runs in all). The variation in the tables is therefore between rules and training samples, not sampling error of the target law; the bootstrap does not cover the uncertainty of training the rule. Uniform randomization over the forms of the rule\textquotesingle s items needs \nRulesUnifLo--\nRulesUnifHi{} times the records of branching (median \nRulesUnifMed), and the ratio within the pairwise model is \nRulesModelLo--\nRulesModelHi{} (median \nRulesModelMed, \nRulesModelN{} runs).}
\newcommand{\nMisspecAppendixText}{Three findings are stable over the \nMisRuns{} runs. First, the forms that are optimal for the pairwise model identify it at the target law: all \nMisStarts{} starts reach one interior maximizer in \nMisWellMF{} runs (several interior maximizers in \nMisMultiMF, none in \nMisNoneMF). On these forms the bias is at most \nMisBiasSdMf\% of the standard deviation of $t$, or \nMisBiasPtsMf{} accuracy points, and the sandwich variance is between \nMisSandMfLo{} and \nMisSandMfHi{} of $\Var_P(t)$, below the variance of the design-based mean on the tree. Second, the other designs identify the model less often. The model-optimal branching design has one interior maximizer from all starts in \nMisWellMB{} runs, several in \nMisMultiMB{} and none in \nMisNoneMB; the optimal design-based forms in \nMisWellF, \nMisMultiF{} and \nMisNoneF; the rule\textquotesingle s own tree in \nMisWellT, \nMisMultiT{} and \nMisNoneT{} (the remaining runs have one interior maximizer that some starts miss). Where both model-optimal designs have an interior maximizer (\nMisBoth{} runs), the ratio of the sandwich variances of forms and branching is \nMisVarMmLo--\nMisVarMmHi, the branching design, optimal for the fitted law, has the larger variance at the true law in \nMisMbWorse{} runs, and with the bias included the mean squared error of the forms relative to the branching design is a median \nMisMseMma, \nMisMseMmb{} and \nMisMseMmc{} at 1,000, 4,000 and 16,000 records (ranges \nMisMseMmLoa--\nMisMseMmHia{} and \nMisMseMmLoc--\nMisMseMmHic{} at the two ends). Where no interior maximizer exists, pairwise maximum likelihood is defined only through its algorithm: in the simulated audits it is EM from the uniform law, stopped when the estimate of the contrast changes by less than $10^{-5}$ or after 500 iterations, and its accuracy there (Table~\ref{tab:curve}) is that of this early-stopped estimate. Third, the variance computed from the model\textquotesingle s information matrix is between \nMisSandModelLo{} and \nMisSandModelHi{} times too small at the interior maximizers, which contributes to the under-coverage of model-based intervals in Table~\ref{tab:em} beyond the effect of bias. The design-based ratio at the same laws is \nMisDesignLo--\nMisDesignHi{} (median \nMisDesignMed).}
